% The article for pull request #68: the Lean ladder's remaining phases.
\documentclass[11pt]{article}
\usepackage[T1]{fontenc}
\usepackage{lmodern}
\usepackage{microtype}
\usepackage[margin=1in]{geometry}
\usepackage{amsmath,amssymb,amsthm,mathtools}
\usepackage{algorithm}
\usepackage{algpseudocode}
\usepackage{booktabs}
\usepackage{array}
\usepackage{enumitem}
\setlist{noitemsep,topsep=2pt}
\usepackage[numbers,sort&compress]{natbib}
\usepackage[hidelinks]{hyperref}
\usepackage[capitalise,noabbrev]{cleveref}
\makeatletter
\renewcommand{\theHALG@line}{\thealgorithm.\arabic{ALG@line}}
\makeatother
\renewcommand{\bibfont}{\footnotesize}

\theoremstyle{plain}
\newtheorem{theorem}{Theorem}[section]
\newtheorem{lemma}[theorem]{Lemma}
\newtheorem{proposition}[theorem]{Proposition}
\newtheorem{corollary}[theorem]{Corollary}
\theoremstyle{definition}
\newtheorem{definition}[theorem]{Definition}
\theoremstyle{remark}
\newtheorem{remark}[theorem]{Remark}

\DeclareUrlCommand\lean{\urlstyle{tt}}
\newcommand{\mlean}[1]{\text{\texttt{#1}}}
\newcommand{\rn}{\operatorname{rn}}
\newcommand{\rnf}{\operatorname{rnf}}
\newcommand{\cosine}{\operatorname{cos}}
\newcommand{\norm}[1]{\lVert #1 \rVert}
\newcommand{\abs}[1]{\lvert #1 \rvert}
\newcommand{\inner}[2]{\langle #1, #2 \rangle}
\newcommand{\cnt}{\operatorname{count}}
\newcommand{\countSV}{\operatorname{count}^{\mathrm{sv}}}
\newcommand{\betaS}{\beta_{\mathrm{store}}}
\newcommand{\betaQ}{\beta_{\mathrm{query}}}
\newcommand{\A}{\mathbb{A}}
\setlength{\emergencystretch}{1.5em}
\renewcommand{\topfraction}{0.95}
\renewcommand{\bottomfraction}{0.7}
\renewcommand{\textfraction}{0.05}
\renewcommand{\floatpagefraction}{0.9}
\setcounter{topnumber}{3}
\setcounter{totalnumber}{4}

\title{Store Models, Float Bounds and Design Theorems:\\
Finishing a Lean Ladder for a Retrieval Algebra}
\author{mind-palace contributors}
\date{November 7, 2026}

\begin{document}
\maketitle

\begin{abstract}
A knowledge store whose retrieval algebra is written to a ledger of laws, each stated in Lean and
attacked before it is proved, reaches a stage at which the remaining claims are of three kinds:
correctness models of the store's own write paths, bounds on what finite arithmetic does to a
ranked retrieval, and theorems a design note stated in prose before any code relied on them. This
article reports the branch that closes those three phases. The correctness models cover a scoped
re-key that keeps every row the store's laws need, with the unscoped migration as a refuted
control; version lineage, under which a path's chain members are exactly its ledger versions; the
crash points of a landing and of a relabel, with cache additions written before the transaction so
that a crash leaves only excess that no search reads; and an existence oracle shown to reveal
exactly the table. The float claims enter the ladder's numerical rung: a rounding model with the
format's precision and exponent range as parameters, attacked bit for bit against the machine's
float32 and the store's float16 write path; a bound of about $2^{-11}$ on what float16 storage
can move a cosine; and a top-$k$ margin certificate, proved and computed on request. Weyl's
inequality in counting form, proved from a dimension count because the library lacks it, gives
the condition under which three tolerances decide zero against nonzero. The design theorems write
a temporal transport as a superconnection and correct two statements of their source. After
integration the gates are 3,131 tests, 281 statements judged identical to their proofs by a
comparator, 46,085 declarations replayed by Lean's kernel and five independent checkers, 90
ledger rows in 13 families, 1,048 theorems and no \texttt{sorry}. A replay of 172,560 ranked slots over an archived
store found 51 changed by float16 storage, every one inside the bound and none certified.
\end{abstract}

\section{Introduction}

The repository under study \citep{mindpalace2026} keeps a knowledge store whose retrieval
algebra is governed by a ledger of laws. Each law has one row: its statement, the structure it
instantiates, the code it anchors, a label for its epistemic status, a Python test that would
catch it false, and the Lean declarations that state it. The Lean side is a ladder: statements as
specifications (L1), proofs over exact arithmetic (L2), and a numerical rung (L3) entered one
claim at a time, in which a floating-point claim is proved over a model of rounding that is itself
attacked against the machine. At the start of this branch the ledger held 87 rows in 12 families
and 248 statements judged by a comparator against their proofs.

At that stage the claims that remain are not the algebraic ones. They are of three kinds. First,
the store's own write paths carry informal correctness arguments: a migration that re-keys rows,
a landing that writes atoms, a transaction and a cache in a fixed order, a relabel that moves a
class, an oracle the write path consults to skip work. Each had a prose argument and a test, and
in two cases the argument was wrong: the migration dropped rows it never read, and a crash between
a landing's transaction and its cache step hid rows from every reader. A model of the store, not
of the mathematics it computes, is what settles these. Second, the store holds its vectors in
binary16, and the search engine scores a query against that column in float32 after rounding the
query to binary16 as well, so every ranked read is ranked by scores that are not the exact cosines
between the query and the embedder's vectors, and nothing said by how much; three tolerances
elsewhere decide whether a computed spectral value is zero, and nothing said when they decide
correctly. Third, a design note on temporal retrieval stated six results in prose, four of them
already machine-checked; the remaining two write the transport of cochains along a chain of times
as a superconnection and ask when it is flat, and what a transport that commutes with the
coboundary owes to its adjoint. Both turned out to be wrong as written.

\paragraph{Contributions.} The branch, 51 commits stacked on \texttt{main}, contributes:
\begin{enumerate}
  \item A model of the store around its shared plane, under which the scoped re-key keeps the
  plane's state, every row outside the note lane, and every row's document, text and lane
  (PLN-6, \lean{Palace.Plane.Store.rekey_keeps}); the unscoped migration drops shed atoms,
  history and edges and merges rows, by theorem and by witness.
  \item The crash points of a landing and of a relabel (PLN-4 extended): with the cache's
  additions written before the transaction and taken back on refusal, a crash leaves the cache
  complete, and two complete planes differing only in their caches answer every search alike
  after any trace (\lean{Palace.Plane.search_after_crash},
  \lean{Palace.Plane.search_after_relabelCrash}).
  \item A model of version lineage under which a path's chain members are exactly its ledger
  versions (\lean{Palace.Lineage.History.mem_chain_iff}), correcting a design note, with a
  witness that a chain's consecutive pairs need not be any commit's diff.
  \item An exact account of the write path's existence oracle: it reveals the table and nothing
  else (\lean{Palace.Plane.known_eq_known_iff}), it interferes with an uncleared write, and
  bounded to a reader's own classes it reveals nothing beyond the view
  (\lean{Palace.Plane.knownIn_congr}).
  \item A rounding model $\rnf\,p\,e_{\min}$ with the format's precision and least exponent as
  parameters, equal to the existing binary64 model at its parameters and attacked against the
  machine's binary32 conversion and a table of the store's binary16 write path; the float16
  cosine bound $\betaS$ and the top-$k$ margin certificate
  (\lean{Palace.Cosine.top_unchanged}), computed on request and replayed over an archived store
  of 17,256 rows (ledger family XIII, VEC-1).
  \item Weyl's inequality in counting form for symmetric operators and for singular values,
  proved from a dimension count, and the spectral-gap lemma that says when a cut keeps exactly
  the kernel below it (SPC-3, \lean{Palace.Gap.count_eq_of_gap},
  \lean{Palace.Gap.countSV_eq_of_gap}); a pseudo-inverse cut moved from the library default
  to one the lemma admits.
  \item The superconnection $\A = d + \tau$ on the nerve of a chain of times, with
  $\A^2 = [d, \tau] + \tau^2$, flatness characterized, and two errata to the source note (TRN-10
  to TRN-15).
  \item Code fixes, each pinned by a test that fails on the old code: a half delete that stranded
  a note, a job claim made a compare-and-set, refusals beside a live writer, a memoized closure on
  a frozen type, algebraic connectivity on disconnected graphs, a row-length roundoff tolerance,
  and the pseudo-inverse cut.
\end{enumerate}

\section{Background and related work}

\paragraph{The ledger and the ladder.} The ledger is a Markdown table, one row per law,
cross-referenced two ways with the Lean package: each cited declaration carries its law id in its
docstring, and a script checks that every marked declaration has a row and every row's citation
exists. Each family has a model (\texttt{Palace/X/Defs.lean}) mirroring the code's definitions,
statements (\texttt{Spec/X.lean}) with proofs left as \texttt{sorry}, and proofs
(\texttt{Palace/X.lean}). A comparator judges the proofs against the statements: each theorem in
\texttt{Spec} is proved in \texttt{Palace} with exactly the same statement from the three standard
axioms, and everything the statements reach is replayed through Lean's kernel and five independent
checkers. A proof file therefore cannot change what the ledger states. Design theorems, proved
before any code relies on them, live in the proof files only and are replayed but not judged.

Each family is attacked before it is proved: property-based tests in the manner of QuickCheck
\citep{claessen2000}, run by Plausible, the Lean~4 library for that method, with a fixed seed on
an executable shadow of the law (the law over $\mathbb{Q}$ where the code works over
$\mathbb{R}$, a decidable transcription where the code's objects are infinite). Every attack has a
control, a nearby false law that must be refuted, so that an attack that survives because it
tests nothing is caught. The toolchain is Lean~4 \citep{demoura2021} and Mathlib
\citep{mathlib2020}.

\paragraph{Floating point and perturbation theory.} The numerical rung states float claims over
a model of IEEE~754 rounding \citep{ieee754}: round to nearest, ties to even, on the grid of a
given precision and least exponent, without overflow. The error bound used,
$2^{-p}\max(\abs{x}, 2^{e_{\min}+p-1})$, is the standard one \citep{goldberg1991,higham2002}; the
summation bound $\gamma_n = nu/(1 - nu)$ that one code fix adopts is Higham's. Backward stability
of the LAPACK decompositions NumPy calls, with a modestly growing polynomial in the dimension, is
as the LAPACK guide states it \citep{lapack1999}. Weyl's inequality for the eigenvalues of a
perturbed symmetric operator goes back to \citet{weyl1912}; its singular-value form is Mirsky's
\citep{mirsky1960}; both are presented with the Courant--Fischer characterization in
\citet{hornjohnson2013} and \citet{stewartsun1990}. The pseudo-inverse is Penrose's
\citep{penrose1955} and its perturbation bound Wedin's \citep{wedin1973}. Effective resistance
through the Laplacian's pseudo-inverse and Rayleigh's monotonicity law are in
\citet{doylesnell1984}; algebraic connectivity is Fiedler's \citep{fiedler1973}; the Hodge
Laplacian of a graph is set out by \citet{lim2020}, and Forman's curvature \citep{forman2003} is
the ledger's curvature family, already proved and not touched here.

\paragraph{Superconnections, noninterference, side channels.} Quillen introduced
superconnections as odd operators $\A = d + \tau$ whose square plays the role of a curvature
\citep{quillen1985}; the design note under study uses the form for the transport of cochains
along a chain of times, and the theorems here are about that discrete instance. The shared plane's
reader law is noninterference in the sense of \citet{goguenmeseguer1982}, proved by Rushby's
unwinding \citep{rushby1992} on an earlier branch under four hypotheses. The existence oracle
characterized here is the deduplication side channel of \citet{harnik2010}: a write path that
skips work for content it already holds tells the writer that someone holds it.

\section{Theory}

The Lean declaration that checks a statement is named after it. Families with a ledger row
(Cosine, Gap, Plane) have their statements in \texttt{Spec/} and are judged by the comparator;
Lineage and Superconnection are design theorems, replayed by the kernel checkers but stated in the
proof files.

\subsection{Rounding to a binary format, and the float16 cosine bound (VEC-1)}

\begin{definition}[Rounding to a format]
Let $K$ be an ordered field with a floor function. For precision $p \ge 1$ and least exponent
$e_{\min} \in \mathbb{Z}$, the exponent of $x$'s unit in the last place is
$c(x) = \max(\lfloor \log_2 \abs{x} \rfloor - (p - 1),\, e_{\min})$, and
\[
  \rnf\,p\,e_{\min}\,(x) = \operatorname{rne}\bigl(x / 2^{c(x)}\bigr)\cdot 2^{c(x)},
\]
with $\operatorname{rne}$ rounding to the nearest integer, ties to even
(\lean{Palace.Format.rnf}). Binary16 is $\rn_{16} = \rnf\,11\,(-24)$ and binary32 is
$\rn_{32} = \rnf\,24\,(-149)$. The existing binary64 model satisfies
$\rn = \rnf\,53\,(-1074)$ (\lean{Palace.Format.rn_eq_rnf}).
\end{definition}

\begin{theorem}[Rounding error]\label{thm:rnf}
For every $x$ (\lean{Palace.Format.abs_rnf_sub_le}),
\[
  \abs{\rnf\,p\,e_{\min}\,(x) - x} \le 2^{-p}\max\bigl(\abs{x}, 2^{e_{\min}+p-1}\bigr).
\]
In particular $\abs{\rn_{16}(x) - x} \le u_{16}\max(\abs{x}, \tau_{16})$ with $u_{16} = 2^{-11}$
and $\tau_{16} = 2^{-14}$ (\lean{Palace.Cosine.abs_rn16_sub_le}).
\end{theorem}

\begin{proof}[Proof sketch]
Rounding moves $x$ by at most half a unit in the last place,
$2^{c(x)}/2$ (\lean{abs_rnf_sub_le_half_ulp}). In the normal range $c(x) = \lfloor \log_2
\abs{x} \rfloor - p + 1$ and $2^{\lfloor \log_2 \abs{x} \rfloor} \le \abs{x}$; below it
$c(x) = e_{\min}$ and the bound is the constant $2^{e_{\min}-1}$.
\end{proof}

\begin{definition}[Vectors and scores]
For $x \in \mathbb{R}^d$ let $\operatorname{round}_{16}(x)$ apply $\rn_{16}$ coordinatewise. The
cosine of $q$ and $x$ is $\cosine(q, x) = \inner{q}{x}/(\norm{q}\norm{x})$, taken as $0$ when
either vector is zero. Set
\[
  \eta(d, \nu) = u_{16}\sqrt{1 + d\,\tau_{16}^2/\nu^2},\qquad
  \betaS(d, \nu) = \frac{2\eta}{2 - \eta},\qquad
  \betaQ(q, q') = \frac{2\norm{q' - q}}{\norm{q} + \norm{q'}}.
\]
A finite set $T$ of indices is a top set of scores $t$ when every member scores at least every
non-member (\lean{Palace.Cosine.IsTop}).
\end{definition}

\begin{theorem}[Storage and query bounds]\label{thm:store}
\begin{enumerate}
  \item $\norm{\operatorname{round}_{16}(x) - x} \le u_{16}\sqrt{\norm{x}^2 + d\,\tau_{16}^2}$
  (\lean{Palace.Cosine.norm_round16_sub_le}).
  \item For any $q, x, y$ in a real inner product space,
  $\abs{\cosine(q, x) - \cosine(q, y)} \le 2\norm{x - y}/(\norm{x} + \norm{y})$
  (\lean{Palace.Cosine.abs_cosine_sub_le}).
  \item If $0 < \nu \le \norm{x}$ and $\eta(d, \nu) < 2$, then
  $\abs{\cosine(q, \operatorname{round}_{16}(x)) - \cosine(q, x)} \le \betaS(d, \nu)$
  (\lean{Palace.Cosine.abs_cosine_round16_sub_le}).
  \item $\abs{\cosine(q', x) - \cosine(q, x)} \le \betaQ(q, q')$
  (\lean{Palace.Cosine.abs_cosine_sub_query}).
\end{enumerate}
\end{theorem}

\begin{proof}[Proof sketch]
(1) sums the squared coordinate bounds, using $\max(\abs{x_i}, \tau_{16})^2 \le x_i^2 +
\tau_{16}^2$. (2) is the inner-product form of the Dunkl--Williams inequality
\citep{dunklwilliams1964}: with $a = \norm{x}$, $b = \norm{y}$, $s = \inner{x}{y}$, the
difference of cosines is $\inner{q}{b x - a y}/(\norm{q} a b)$, Cauchy--Schwarz bounds the
numerator, and $\norm{bx - ay}^2 (a + b)^2 \le 4a^2b^2\norm{x - y}^2$ follows from
$\norm{bx - ay}^2 = 2ab(ab - s)$ and $\norm{x - y}^2 = a^2 - 2s + b^2$. (3) applies (2) to $x$
and its rounding, which by (1) moves it by $E \le \eta\norm{x}$, so that
$2E/(2\norm{x} - E) \le 2\eta/(2 - \eta)$. (4) is (2) with query and vector swapped.
\end{proof}

\begin{theorem}[Margin certificate]\label{thm:margin}
Let $s, t : \kappa \to \mathbb{R}$ with $\abs{s_j - t_j} \le \beta$ for every $j$, and let $S$ be a
finite set such that $s_j + 2\beta < s_i$ for every $i \in S$ and $j \notin S$. Then
$t_j < t_i$ for every such pair, and every top set of $t$ with $\abs{S}$ elements equals $S$
(\lean{Palace.Cosine.top_eq_of_margin}).
\end{theorem}

\begin{theorem}[VEC-1, the certificate the store computes]\label{thm:vec1}
Let each stored vector be $\operatorname{round}_{16}(y_j)$ for a vector $y_j$ with
$\nu \le \norm{y_j}$, let the store score with $q'$ in place of the query $q$, and let each
computed score satisfy $\abs{s_j - \cosine(q', \operatorname{round}_{16}(y_j))} \le \gamma$. If
the members of $S$ beat every non-member by more than $2(\betaS(d, \nu) + \betaQ(q, q') +
\gamma)$, then the exact cosines $\cosine(q, y_j)$ rank $S$ strictly above the rest, and $S$ is
their only top set of its size (\lean{Palace.Cosine.top_unchanged}).
\end{theorem}

\begin{proof}[Proof sketch]
\Cref{thm:margin}: if $T \ne S$ has the same size, some $i \in S \setminus T$ and $j \in T
\setminus S$ exist, and $t_j < t_i$ contradicts $T$ being a top set. \Cref{thm:vec1} chains the
three bounds of \cref{thm:store} through the triangle inequality, so each score is within
$\beta = \betaS + \betaQ + \gamma$ of its exact cosine.
\end{proof}

The ledger row fixes $\nu = 1/2$ as a hypothesis on the corpus (the embedder returns unit vectors;
stored norms measured 0.99988 to 1.00012) and $\gamma = 2(d + 4)\,2^{-24}$ as the stated bound on
the float32 distance kernel (the classical bound for a dot product and two norms of $d$ terms, a
square root and a division; measured within $5.6\times 10^{-8}$); at $d = 768$,
$\betaS = 4.884\times 10^{-4}$.

\subsection{Weyl's inequality counted, and the gap lemma (SPC-3)}

\begin{definition}
Let $E$ be a real inner product space of finite dimension $n$ and $T$ a symmetric operator on
it with eigenvalues $\lambda_1, \dots, \lambda_n$ (Mathlib's \lean{IsSymmetric.eigenvalues}).
$\cnt_T(t)$ is the number of $i$ with $\lambda_i \le t$ (\lean{Palace.Gap.count}). $B$ is
within $\varepsilon$ of $A$, written $\operatorname{Near}(A, B, \varepsilon)$, when
$\abs{\inner{x}{(B - A)x}} \le \varepsilon\norm{x}^2$ for every $x$. For any operator $M$,
$\countSV_M(t) = \cnt_{M^\dagger M}(t^2)$ counts the singular values at most $t$
(\lean{Palace.Gap.countSV}).
\end{definition}

\begin{lemma}[The dimension count]\label{lem:dim}
If fewer than $n$ eigenvalues of $T_1$ exceed $c_1$ and of $T_2$ are at most $c_2$, counted
together, some $x \ne 0$ satisfies $\inner{x}{T_1 x} \le c_1\norm{x}^2$ and
$c_2\norm{x}^2 < \inner{x}{T_2 x}$ (\lean{Palace.Gap.exists_of_card_lt}).
\end{lemma}

\begin{proof}[Proof sketch]
The map $x \mapsto (\inner{b^1_i}{x})_{i \in S_1} \times (\inner{b^2_i}{x})_{i \in S_2}$ into a
space of dimension $\abs{S_1} + \abs{S_2} < n$ has a nonzero kernel element $x$. Expanding each
quadratic form in its eigenbasis, the terms that could exceed $c_1$ (respectively fall below
$c_2$) vanish, and Parseval gives the two inequalities; the second is strict because some
coefficient of $x$ in the $T_2$ basis is nonzero.
\end{proof}

\begin{theorem}[Weyl counted]\label{thm:weyl}
If $\operatorname{Near}(A, B, \varepsilon)$ for symmetric $A, B$, then for every $t$,
$\cnt_A(t - \varepsilon) \le \cnt_B(t) \le \cnt_A(t + \varepsilon)$
(\lean{Palace.Gap.count_sub_le}, \lean{Palace.Gap.count_le_add}). If
$\norm{(M' - M)x} \le \varepsilon\norm{x}$ for every $x$, then $\countSV_M(t - \varepsilon) \le
\countSV_{M'}(t)$ for $\varepsilon \le t$ and $\countSV_{M'}(t) \le \countSV_M(t + \varepsilon)$
for $0 \le \varepsilon$, $0 \le t$ (\lean{Palace.Gap.countSV_sub_le},
\lean{Palace.Gap.countSV_le_add}).
\end{theorem}

\begin{proof}[Proof sketch]
Suppose $\cnt_A(t - \varepsilon) > \cnt_B(t)$. Since the eigenvalues above $t - \varepsilon$
and those at most $t - \varepsilon$ are all $n$, \cref{lem:dim} with $c_1 = t - \varepsilon$,
$c_2 = t$ gives $x$ with $\inner{x}{Ax} \le (t - \varepsilon)\norm{x}^2 < \inner{x}{Bx} -
\varepsilon\norm{x}^2$, against nearness. The singular-value forms use $\inner{x}{M^\dagger Mx}
= \norm{Mx}^2$ and the triangle inequality $\norm{M'x} \le \norm{Mx} + \varepsilon\norm{x}$.
\end{proof}

\begin{theorem}[The gap lemma]\label{thm:gap}
Let every nonzero eigenvalue of symmetric $A$ be at least $g$, let
$\operatorname{Near}(A, B, \varepsilon)$, and let $0 \le \varepsilon \le \theta$ with
$\theta + \varepsilon < g$. Then $\cnt_B(\theta) = \cnt_A(0)$: the cut $\theta$ keeps exactly
the kernel's dimension below it (\lean{Palace.Gap.count_eq_of_gap}). The same holds for
singular values, with $g^2$ a lower bound on the nonzero eigenvalues of $M^\dagger M$
(\lean{Palace.Gap.countSV_eq_of_gap}).
\end{theorem}

\begin{proof}[Proof sketch]
$\cnt_A(\theta + \varepsilon) = \cnt_A(0)$ because no eigenvalue lies in $(0, g)$, and
$\cnt_A(0) \le \cnt_A(\theta - \varepsilon)$ because $\theta - \varepsilon \ge 0$. Squeeze
with \cref{thm:weyl}.
\end{proof}

The ledger row instantiates $\varepsilon$ from backward stability: $8nu\norm{A}$ for the
decomposition, with $u = 2^{-53}$, plus $nu\norm{L}$ for forming a weighted Laplacian's degree
sums in float64, so $\varepsilon = 9nu\norm{L}$ for the pseudo-inverse cut and $8nu\norm{L_1}$
for the Hodge rank cut, whose integer entries are formed exactly. With
$\theta = 10^{-10}\sigma_{\max}$ the first condition holds up to about $1.0\times 10^5$ notes
($1.1\times 10^5$ edges for the Hodge cut); the second is the gap condition
$g > (10^{-10} + 16nu)\norm{A}$, checked on fixtures. The rise threshold of the reconnection
scan, $10^{-9}$ on a change of conductance, rests further on Wedin's bound \citep{wedin1973},
cited and not formalized: a computed conductance $\hat c$ is within
$\delta = c\,\hat c\,\bigl(4\mu\varepsilon/(g(g - \varepsilon)) + 4(n + 3)u/(g - \varepsilon)\bigr) +
u\hat c$ of the exact $c$, $\mu = (1 + \sqrt 5)/2$, and an exact change of zero reads as no rise
when $2\delta \le 10^{-9}$.

\subsection{The plane's write paths (PLN-4, PLN-6) and the oracle}

The plane model is a state $s$ with, per class $c$, a relation $s.\mathrm{occ}\,c$ of occupancy
rows (path, digest, chunk index, atom, currency flag) and a supersession log; a table of atoms;
and a cache giving each atom two sets of classes, $\mathrm{held}$ and $\mathrm{ever}$. The cache
is \emph{coherent} when it equals, for each table atom, the classes whose rows hold the atom, and
\emph{complete on $C$} when it contains those classes for $c \in C$
(\lean{Palace.Plane.Coherent}, \lean{Palace.Plane.Complete}); $\operatorname{WF}$ says no row
names an atom the table lacks. A search under clearance $C$ prefilters by the cache and keeps only
atoms a cleared relation holds, so a cache entry in excess is never returned.

\begin{definition}[The landing's crash points]
For a fiber $f$ (one document version's rows, in class $f.\mathrm{cls}$), the code's order is:
insert the missing atoms; add $f.\mathrm{cls}$ to the cache of every atom of $f$ that the table
holds (\lean{addCache}); the transaction (\lean{commit}); then refresh the cache of every affected
atom to its truth (\lean{refresh}). $\mlean{landAdd}\,f\,s$ is the state after the additions and
before the transaction; $\mlean{landAddCommit}\,f\,s$ the state after the transaction (or after
its rollback, when the land is refused) and before the refresh. $\mlean{land}\,f\,s$ is the whole
write. $\operatorname{SameRel}(s, t)$ says $s$ and $t$ agree on relations, logs and table and
differ at most in their caches. A relabel from class $o$ to class $n$ has the same shape:
$\mlean{addRelabel}$ adds $n$ to the cache of every atom $o$'s rows hold, $\mlean{relabelCommit}$
moves the rows and the log, and a refresh removes $o$.
\end{definition}

\begin{theorem}[PLN-4, crash points]\label{thm:crash}
From a plane complete on every class with no dangling row:
\begin{enumerate}
  \item $\mlean{landAdd}\,f\,s$ and $\mlean{landAddCommit}\,f\,s$ are complete on every class
  (\lean{Palace.Plane.complete_landAdd}, \lean{complete_landAddCommit}).
  \item When the land is not refused, refreshing after $\mlean{landAddCommit}$ gives exactly
  $\mlean{land}\,f\,s$; when it is refused and the additions are taken back, the result is again
  $\mlean{land}\,f\,s$, from a coherent plane (\lean{land_of_landAdd},
  \lean{land_of_landAdd_refused}).
  \item Two planes with $\operatorname{SameRel}$, both complete with no dangling row, answer
  every search alike after any trace of writes (\lean{search_run_of_sameRel}); hence from a
  coherent plane, every search after any trace following $\mlean{landAddCommit}\,f\,s$ equals
  the one following $\mlean{land}\,f\,s$ (\lean{search_after_crash}).
  \item The relabel analogues (\lean{complete_addRelabel}, \lean{complete_relabelAdd},
  \lean{relabel_of_addRelabel}, \lean{relabel_of_addRelabel_clash},
  \lean{coherent_relabel}, \lean{search_after_relabelCrash}).
\end{enumerate}
\end{theorem}

\begin{proof}[Proof sketch]
Additions can only enlarge a cache entry, so completeness survives them; the transaction gives
rows to atoms that already name the row's class. The refresh rewrites a touched atom's whole
cache entry to its truth, which is why both orders end in the same state. For (3), every write
preserves $\operatorname{SameRel}$ and completeness, the view under $C$ depends on relations
alone, and a search keeps only what a cleared relation holds, so a cache in excess changes no
answer.
\end{proof}

Two witnesses bound what the theorem does not say. Re-running a land after a crash does not
restore coherence: the first run's freed atoms stay cached in excess, complete all the same, until
\lean{repair} (\lean{Palace.Plane.Witness.rerun_keeps_excess}). With the cache written last, a
relabel that crashed after its transaction hid the moved atoms from the new class, and a re-run
refreshed nothing because the old relation was empty
(\lean{Palace.Plane.Witness.relabel_rerun_hides}).

\begin{definition}[The store around the plane]
A physical row (\lean{Palace.Plane.Phys}) has an id, an optional document, a text, and a flag
saying whether its class is a plane class; a \emph{note record} is a plane-class row with a
document. A store is a plane and a finite set of physical rows. The scoped re-key
$\mlean{Store.rekey}\,\mathrm{key}$ gives each note record the id $\mathrm{key}(\mathrm{path},
\mathrm{text})$ and touches nothing else; the unscoped $\mlean{Store.rekeyReset}$ reads the
documents' rows, drops the plane, and lands the current rows again from the empty plane.
\end{definition}

\begin{theorem}[PLN-6]\label{thm:rekey}
For every key and store $s$: the re-key leaves the plane's state unchanged, keeps every row that
is not a note record, gives every note record its key, and keeps the set of (document, text, lane)
contents (\lean{Palace.Plane.Store.rekey_keeps}); from a $\operatorname{Good}$ plane
(coherent, no dangling row, every path at one stratum) it returns a $\operatorname{Good}$ plane
with the same view under every clearance (\lean{Store.rekey_good}); and it is idempotent
(\lean{Store.rekey_idem}). The unscoped form leaves no row without a document
(\lean{Store.rekeyReset_drops_atoms}); on a store with one edited note and four code rows it
drops the first version and its edge (\lean{Witness.rekeyReset_drops_history}) and leaves one of
the four rows, two shed atoms unread and two records merged (\lean{Witness.rekeyReset_merges}).
\end{theorem}

\begin{definition}[The oracle]
$\mlean{known}\,s\,S$ is the set of named atoms $S$ the table holds; $\mlean{knownIn}\,C\,s\,S$ is
those among them that some relation of $C$ references.
\end{definition}

\begin{theorem}[What the oracle reveals]\label{thm:known}
$\mlean{known}\,s = \mlean{known}\,t$ if and only if $s$ and $t$ have the same table
(\lean{Palace.Plane.known_eq_known_iff}); beyond a reader's view it gives, for each named atom no
cleared relation references, one bit, whether an uncleared row holds the atom or the atom is
dormant (\lean{known_eq}, \lean{mem_known_iff}); an uncleared land changes it while the view stays
(\lean{Witness.known_interferes}); and two planes with the same view under $C$ give the same
$\mlean{knownIn}\,C$ (\lean{knownIn_congr}).
\end{theorem}

\subsection{Lineage}

\begin{definition}
A history gives each commit an optional first parent and a tree, the blob at each path. The
ledger is a list $L$ of commits. The rows at path $p$ are the pairs (blob before, blob after) of
the commits in $L$ whose blob at $p$ differs from the first parent's (\lean{History.rows}); the
chain threads them: a row's old blob starts an empty chain, and its new blob is appended unless
the chain already ends there (\lean{Palace.Lineage.step}, \lean{thread}). A version at $p$ is a
blob some ledger commit holds there.
\end{definition}

\begin{theorem}\label{thm:lineage}
If $L$ holds every first parent of its commits and a rank $r$ decreases along first parents, then
a blob is a chain member at $p$ if and only if it is a version at $p$
(\lean{Palace.Lineage.History.mem_chain_iff}). On a history with one side branch of five commits
the chain is every version in the ledger's order, HEAD's first-parent line holds three of the
five, and three consecutive chain pairs are no commit's diff
(\lean{Palace.Lineage.Witness.side_chain}).
\end{theorem}

\begin{proof}[Proof sketch]
A chain member is the old or new blob of some row, so some ledger commit or its first parent
holds it. Conversely, walk a version's first-parent ancestry by strong induction on $r$ to the
commit whose blob differs from its parent's, where a row introduced it
(\lean{History.new_of_version}). The witness is decided by computation.
\end{proof}

\subsection{The superconnection (TRN-10 to TRN-15)}

A citation complex at a time is a type of notes $A$ with a relation $X$; a cochain is a function
on notes and a function on ordered pairs, read on the citations through a mask
(\lean{Palace.Transport.Cochain}). The coboundary $d_X$ sends a note cochain $\varphi$ to the drop
$\varphi(v) - \varphi(u)$ on each citation and an edge cochain to zero (\lean{dC}). The transport
$\tau_\sigma$ along a note map $\sigma : A \to B$ pulls a note cochain back by $\sigma$ and an edge
cochain through the target's citations, reading $0$ on a citation whose image is not a citation
(\lean{tC}, on \lean{pull} and \lean{pullEdge}). On a preorder of times $T$, the total complex
$\operatorname{Tot}$ holds a cochain for each time, each pair $s < t$ and each triangle
$s < t < u$, valued at the earliest time: the nerve to its 2-skeleton. $d$ applies each time's
coboundary with the sign $(-1)^q$ of the time degree (\lean{dTot}); $\tau$ is the \v{C}ech
differential with the transports as coefficients, $(\tau c)_{st} = \Theta_{st} c_t - c_s$ and
$(\tau e)_{stu} = \Theta_{st} e_{tu} - e_{su} + e_{st}$ (\lean{cech}).

\begin{theorem}[TRN-10, TRN-11]
In any ring with $d^2 = 0$, $(d + \tau)^2 = (d\tau + \tau d) + \tau^2$
(\lean{Palace.Transport.superconnection_sq}). On $\operatorname{Tot}$, $d^2 = 0$
(\lean{dTot_mul_self}) and $\operatorname{face}^2 = 0$ (\lean{face_mul_self}); the bracket on
a pair $s < t$ is minus TRN-5's curvature of $\sigma_{st}$ at the note cochain of $t$, read on the
citations of $s$ (\lean{bracket_trans_apply}, \lean{dC_tC_sub}); and $\tau^2$ on a triangle is
the composition defect $\Theta_{st}\Theta_{tu} - \Theta_{su}$ (\lean{cech_mul_self_apply}).
\end{theorem}

\begin{theorem}[TRN-12, TRN-13: flatness]\label{thm:flat}
With transports $\tau_\sigma$ along note maps $\sigma_e$ on every pair of times:
\begin{enumerate}
  \item $\tau^2 = 0$ if and only if on every triangle $s < t < u$ the notes compose,
  $\sigma_{tu}\circ\sigma_{st} = \sigma_{su}$, and no citation is restored: a citation of $s$
  whose image under $\sigma_{su}$ is a citation of $u$ has an image under $\sigma_{st}$ that is a
  citation of $t$ (\lean{tau_sq_eq_zero_iff}). Composition alone does not suffice: with the
  identity on notes, a citation present at times 0 and 2 and severed at time 1 gives
  $\tau^2 \ne 0$ (\lean{tau_sq_ne_zero_of_restored}).
  \item $\A = d + \tau$ is flat if and only if every transport is a chain map (F1: each citation
  of the source collapses or carries forward) and every triangle composes with no restoration
  (\lean{flat_iff}); flat if and only if a bicomplex (\lean{flat_iff_bicomplex}); on citation
  complexes without loops, flat if and only if F1 and F2, as the source note states
  (\lean{flat_iff_of_loopless}), because F1 rules restoration out (\lean{noRestore_of_F1}).
  \item On a flat $\A$ the two routes around a diamond of times carry every note the same way
  (\lean{diamond_routes_agree}).
\end{enumerate}
\end{theorem}

\begin{theorem}[TRN-14: the correction]\label{thm:corr}
In any ring with $d^2 = 0$ and $[d, h] = [d, \tau]$, $(d + \tau - h)^2 = (\tau - h)^2$ while
$(d + \tau + h)^2 = 2[d, \tau] + (\tau + h)^2$ (\lean{sub_correction_sq},
\lean{add_correction_sq}). On two times with one citation severed and its ends still joined,
$d + \tau - h$ is flat and $d + \tau + h$ is not (\lean{correction_sign}). Flatness of
$d + \tau - h$ needs the corrected transports to compose on every triangle, which the class does
not give (\lean{corrections_not_compose}). A correction that is zero on note cochains and has
the transport's bracket exists if and only if every citation's ends stay in one component of the
target's citations (\lean{exists_correction_iff}); unconstrained, $h = \tau$ always serves
(\lean{flat_sub_first_self}).
\end{theorem}

\begin{theorem}[TRN-15: the adjoint]\label{thm:adj}
The divergence $\operatorname{div}_X\psi(a) = \sum_x (\psi_X(x, a) - \psi_X(a, x))$ is the adjoint
of the coboundary under the counting inner products (\lean{sum_mask_coboundary_mul}). If
$\sigma$ has every fiber of size one, F1 holds, and every citation between images is the image of
a citation, then $\sigma^*\operatorname{div}_{X'} = \operatorname{div}_X\sigma^*$
(\lean{pull_div_eq_div}), and such a transport commutes with the edge Laplacian
(\lean{pullEdge_laplacian}). The inclusion is strict: a citation added at the later time
(\lean{new_citation_not_metric}) and a merge of two notes (\lean{merge_not_metric}) are chain maps
that do not commute with the divergence. A bijection is not needed: winding a six-cycle twice
around a three-cycle merges every note and still commutes (\lean{wind_metric}).
\end{theorem}

\section{Concepts and intuition}

\paragraph{Excess is invisible, absence is not.} The plane's search uses its cache only to choose
candidates and then checks each against a cleared relation. A cache entry that names a class
which does not hold the atom costs a wasted candidate; a missing entry hides an atom from every
reader of that class. So the safe side of any crash is excess, and writing the additions before
the transaction puts every crash on that side: before the transaction the entry is premature,
after it the entry is true and the removals are owed, and neither changes a search. The old order
put the crash on the wrong side, and a re-run did not help because it refreshes only what it
lands and frees. \Cref{thm:crash}(3) makes this precise: two planes that agree on everything but
their caches, both complete, are indistinguishable to every future search.

\paragraph{Why a margin of $2\beta$.} If each computed score is within $\beta$ of its exact
cosine, a member's exact score is at least its computed score less $\beta$ and a non-member's at
most its computed score plus $\beta$, so a gap of more than $2\beta$ between the worst member and
the best non-member survives the move to exact scores. The certificate is about the set, not the
order within it, and the control with a margin of $\beta$ is refuted within nine instances. The
storage bound is a worst case: a vector whose every coordinate sits on a float16 tie moves by a
full $2^{-11}$ relative in each coordinate, the error is orthogonal to the vector, and a query
along the error sees nearly all of it. Typical errors are far smaller, so the certificate settles
only part of the reads, and its refusals land on near-ties, which are exactly the reads float16
can change.

\paragraph{The dimension count, and why a cut needs $\varepsilon \le \theta$.} Mathlib has the
spectral theorem and Rayleigh's quotient for the extreme eigenvalues but neither Weyl's inequality
nor the min-max characterization. The counting forms need less than either: if too few
eigenvectors of $T_1$ lie above $c_1$ and of $T_2$ at most $c_2$ together, a vector orthogonal to
all of them exists and its two quadratic forms disagree by more than the perturbation allows; the
same count with $\norm{Mx}^2$ for the forms gives Mirsky's inequality. A cut $\theta$ below the
backward error can read a kernel value as nonzero, since Weyl allows a zero eigenvalue to be
reported anywhere in $[-\varepsilon, \varepsilon]$, and NumPy's default $10^{-15}\sigma_{\max}$ is
below $8nu\sigma_{\max}$ from $n = 2$ on. The other condition, $\theta + \varepsilon < g$, is the
gap: a nonzero eigenvalue below the cut plus the error is lost.

\paragraph{The sign of the correction.} The bracket $[d, \tau] = d\tau + \tau d$ is the graded
commutator of two odd operators. If $h$ has the same bracket with $d$ as $\tau$ does, subtracting
$h$ removes the bracket from $(d + \tau - h)^2$ and leaves $(\tau - h)^2$; adding it doubles the
bracket, and the source note wrote $d + \tau + h$. The second erratum is of the same kind: the
note's condition for $\tau^2 = 0$ named only composition of the note maps, but $\tau$ carries
edge cochains through the target's citations, and a citation severed at a middle time and present
again later gives a transport that does not factor.

\paragraph{One bit per named atom.} The oracle tells the write path which atoms the table holds,
and \cref{thm:known} says that is all it tells: two planes answer alike if and only if their
tables agree, and bounded to the classes a role reads it adds nothing to the view. The branch
bounds the one caller whose answer reached a reader, a script's embed charge, and leaves the
other, which skips embeds, as it was; on the live store no atom was held by two classes, so the
bound cost no embed.

\section{Implementation}

\subsection{The certificate}

The store computes \cref{thm:vec1}'s certificate on request through \texttt{search\_ranked}
(\cref{alg:cert}); \texttt{search} is unchanged in cost and candidates. Each source is asked for
$k + 1$ candidates so that the best vector outside the top $k$ is present, and the gap is taken
over distinct stored vectors, since two rows of one atom share its vector and its score. The
bound is evaluated in float64 and multiplied by $1 + 2^{-20}$ to cover that evaluation's roundoff.
The query bound is taken against the query rounded to float32 and then to float16, which is what
the engine scores: a query with a coordinate $1 + 2^{-11} + 2^{-30}$, just above a float16 tie,
distinguishes that from rounding once, and the engine's own distance is the witness.

\begin{algorithm}[htbp]
\caption{The margin certificate of a ranked read (\texttt{core/kernel/margin.py}).}\label{alg:cert}
\begin{algorithmic}[1]
\Require query $q \in \mathbb{R}^d$; candidates $(d_1, \kappa_1), \dots, (d_m, \kappa_m)$ with
distances ascending and $\kappa_i$ the stored vector's key; $k \ge 1$
\Ensure a margin (gap, bound) certified when gap $>$ bound
\State $q' \gets \rn_{16}(\rn_{32}(q))$ coordinatewise
\State $\betaQ \gets 2\norm{q' - q}/(\norm{q} + \norm{q'})$ \Comment{$0$ when both norms vanish}
\State $\eta \gets 2^{-11}\sqrt{1 + d\,2^{-28}/\nu^2}$ with $\nu = 1/2$; $\betaS \gets 2\eta/(2-\eta)$
\State $\gamma \gets 2(d + 4)\,2^{-24}$
\State bound $\gets 2(\betaS + \betaQ + \gamma)(1 + 2^{-20})$
\If{$m \le k$} \Return (gap $= \infty$, bound) \EndIf
\State inside $\gets \{\kappa_1, \dots, \kappa_k\}$; worst $\gets d_k$
\State next $\gets$ the first $d_j$ with $j > k$ and $\kappa_j \notin$ inside, else $\infty$
\State \Return (gap $=$ next $-$ worst, bound)
\end{algorithmic}
\end{algorithm}

\subsection{The crash-safe landing order}

\Cref{alg:land} is \texttt{SharedPlane.land} as the branch writes it: presence and cache read in
one scan, additions written only where an atom's cache lacks a class, a rollback taking the
additions back, the removals after the commit. The relabel has the same three steps. The order is
the model's \lean{landAdd}, \lean{landAddCommit} and \lean{refresh}, so \cref{thm:crash} is about
this procedure.

\begin{algorithm}[htbp]
\caption{A landing with the cache's additions before its transaction (\texttt{core/stores/plane.py}).}\label{alg:land}
\begin{algorithmic}[1]
\Require fibers $F$ (each a version's rows in one class), vectors for new atoms
\Ensure the state of \lean{land}, or a refusal that leaves the cache exact
\State wanted $\gets$ the atoms of $F$; gains$[a] \gets$ the classes whose fibers hold $a$
\State take the writer lock; check the geometry
\State cached $\gets$ held(wanted) \Comment{which atoms exist, and their cache}
\State missing $\gets$ wanted $\setminus$ dom(cached); refuse if a missing atom has no vector
\State insert the missing atoms \Comment{D8: atoms first}
\State for each $a \in$ wanted whose cache lacks gains$[a]$: add gains$[a]$ to held$[a]$ and ever$[a]$
\State \textbf{begin} transaction
\For{each fiber $f \in F$}
  \State refuse (rollback, then refresh wanted, then raise) if $f$'s path lives in another class
  \State write $f$'s rows; supersede the path's previous version; log the edge
\EndFor
\State \textbf{commit}
\State refresh the cache of every affected atom to its truth \Comment{the removals, last}
\end{algorithmic}
\end{algorithm}

\subsection{The compare-and-set claim}

The job queue's \texttt{claim} selected a row under a per-process lock and marked it running with
an unguarded update, so two processes on one queue file could both take a job.
\Cref{alg:claim} guards the update with the state and re-selects on a lost race, so the database
decides who holds the row; a \texttt{kinds} filter lets an in-process runner claim only the kinds
it handles, and the conversation runtime, which completed queued rows without claiming them, now
claims through this one site.

\begin{algorithm}[htbp]
\caption{\texttt{JobQueue.claim} as a compare-and-set (\texttt{scheduler/queue.py}).}\label{alg:claim}
\begin{algorithmic}[1]
\Require optional loaded key, blocked tiers, kinds; the clock
\Ensure the claimed job, or none
\Loop
  \State rows $\gets$ \textsc{select} jobs \textsc{where} state $=$ queued, ordered by priority, id
  \State eligible $\gets$ rows not in a blocked tier and, if kinds given, of those kinds
  \If{eligible $= \emptyset$} \Return none \EndIf
  \State chosen $\gets$ the eligible job of least effective priority, the loaded key first, then FIFO
  \State $n \gets$ \textsc{update} jobs \textsc{set} state $=$ running, started, attempts $+ 1$, lease
  \Statex \hspace{\algorithmicindent}\hspace{\algorithmicindent} \textsc{where} id $=$ chosen \textsc{and} state $=$ queued
  \State commit
  \If{$n = 1$} \Return get(chosen) \EndIf \Comment{else another process took it: choose again}
\EndLoop
\end{algorithmic}
\end{algorithm}

\subsection{The other fixes}

\emph{The scoped re-key.} \texttt{rekey\_store} now re-keys the plane's occupancies in one
transaction that updates one column, and a table that predates the plane re-keys its note records
in place, adding rows under their new ids before deleting the old ones so that a crash leaves both
and a re-run finishes; an old id that is another record's key is refused. Nothing else is read or
written.

\emph{The half delete.} A vault delete wrote its tombstone, the delete's key, before retiring the
plane's rows, so a crash between the two left the note current and no rescan returned to it; the
tombstone is written last now. The formal audit then found the state the reversed order opens: a
delete that dies after its retire leaves the key active over retired rows, and if the same bytes
come back before the next rescan an UNCHANGED check that reads the key alone skips them for ever.
UNCHANGED now requires the path's current rows (\texttt{holds\_current}); the ingress writer's
retire had the same state and the same fix.

\emph{Refusals beside a live writer.} The chat ingest verb and the rebuild script, whose first act
drops the table, both ran beside a live daemon with no guard; both refuse now, and the rebuild
holds the supervisor lock for its whole duration, so a daemon started between the check and the
reset is refused in turn. The REPL leaves delegated work to a live daemon.

\emph{Algebraic connectivity.} The normalized Laplacian gives an isolated vertex an identity row,
which puts $1$ in the spectrum, so \texttt{fiedler} read a path plus an isolated vertex as
$\lambda_2 = 0.5$. It now counts components first, on the positive-weight support (the sparse
component count takes a stored zero as an edge), and returns $0$ for a disconnected graph with a
vector separating the first component from the rest.

\emph{Tolerances and the spine.} The dense-Laplacian equivalence test compared degrees against a
fixed $4\epsilon\max$, which the ledger proves only for rows of at most four weights; it now uses
$2\gamma_n\sum\abs{w}$ per row, with a 64-weight fixture and a float32-accumulated control outside
the bound. The pseudo-inverse behind effective resistance passes $\mathrm{rtol} = 10^{-10}$
instead of the library default. \texttt{Spine.restrict}, which builds a stratum's transitive
closure, is memoized per instance on a frozen dataclass, and the real-store clock-law check runs
on generator edges only, sound by CLK-1.

\section{Evaluation}

\subsection{What is proved}

\Cref{tab:gates} gives the gates before and after the branch, from \texttt{formal/check.sh} and
the test suite on the integrated head. The comparator judged all 281 statements identical to
their proofs; the export of 46,085 declarations was accepted by Lean's kernel and by
\texttt{leanchecker-paranoid}, \texttt{lean4lean}, \texttt{nanoda}, \texttt{con-leche} and
\texttt{con-ron}; the axioms behind every declaration are \texttt{propext},
\texttt{Classical.choice} and \texttt{Quot.sound}; the \texttt{sorry} count is zero.

\begin{table}[htbp]
\centering\footnotesize
\caption{The gates before and after the branch, from \texttt{README.md} at the two heads.}\label{tab:gates}
\begin{tabular}{lrr}
\toprule
 & before & after \\
\midrule
ledger rows (families) & 87 (12) & 90 (13) \\
statements judged by the comparator & 248 & 281 \\
cited declarations & 254 & 287 \\
theorems proved in \texttt{formal/Palace} & 912 & 1,048 \\
proof modules / models / statement files & 27 / 26 / 21 & 32 / 31 / 23 \\
lines of models and proofs / statements / attacks & 15,563 / 2,620 / 3,072 & 18,575 / 2,991 / 3,779 \\
attacks / negative controls (modules) & 117 / 104 (23) & 128 / 120 (27) \\
declarations replayed by six checkers & 45,523 & 46,085 \\
deterministic tests & 3,036 & 3,131 \\
\texttt{sorry} & 0 & 0 \\
\bottomrule
\end{tabular}
\end{table}

\Cref{tab:attacks} lists the attacks the branch adds and their controls. An attack survives when
every drawn instance that passes its guard satisfies the law; a control is refuted when Plausible
finds a counterexample.

\begin{table}[htbp]
\centering\footnotesize
\caption{The branch's attacks and controls (\texttt{formal/PalaceTest}, fixed seeds). A control
is a nearby false law; its number is the instances Plausible needed to refute it (for the table
attack, the entries it gets wrong).}\label{tab:attacks}
\setlength{\tabcolsep}{4pt}
\begin{tabular}{p{0.43\textwidth}rp{0.30\textwidth}r}
\toprule
attack & survived & control & refuted \\
\midrule
$\rn_{32}$ against the machine's binary32 (500 doubles) & 500 & ties away; 23 bits; flush to zero & 17 \\
$\rn_{16}\circ\rn_{32}$ against the store's write path (234 doubles) & 234 & round once; ties away; 10 bits; flush & 39; 60; 34; 51 \\
float16 coordinate bound, over $\mathbb{Q}$ & 500 & no floor; half the bound & 7 \\
margin certificate at $2\beta$ & 292 & margin $\beta$ & 9 \\
Weyl's counts on symmetric $3\times 3$ over $\mathbb{Q}$ & 500 & no slack $\varepsilon$ & 9 \\
land crash points leave the cache complete & 500 & cache written last & 6 \\
relabel crash points leave the cache complete & 500 & cache written last & 15 \\
scoped re-key on random physical rows & 500 & the unscoped re-key & 7 \\
(the unscoped re-key's plane) & & re-landing current rows keeps the view & 16 \\
transports compose when notes do, no restoration & 500 & F2 alone & 46 \\
a window with F1 and composition is flat & 500 & middle citations not closed & 13 \\
curvature vanishes on locally constant cochains & 500 & the class without the connection & 15 \\
a bijection carrying exactly the citations commutes with $\delta$ & 500 & a merge; a new citation & 30; 28 \\
\bottomrule
\end{tabular}
\end{table}

\subsection{What is measured}

\paragraph{The float16 replay.} \texttt{eval/float16\_replay.py} copies a vector store archived
before the float16 migration into a temporary directory, builds a float32 and a float16 table from
its 17,256 current code rows at 768 dimensions, and asks each row for its ten nearest other rows
through the search engine's own kernels, once with the float32 row as the query and once with its
float16 rounding. \Cref{tab:replay} reports what float16 storage changed and what the certificate
said. The widest gap among the changed queries is $2.7\times 10^{-5}$, more than forty times
inside the bound of about $1.2\times 10^{-3}$ for a float16 query, and none of them was certified.
The design note's earlier figures (55 changed slots, 63.5\% certified) came from a NumPy reading
and were not reproducible through the engine; the formal audit asked for this replay. It does not
show that float16 never changes a certified read on other corpora; it shows that on this one the
changed reads are the near-ties the certificate refuses.

\begin{table}[htbp]
\centering\footnotesize
\caption{The replay over the archived store: 17,256 queries, 172,560 top-ten slots. The last two
rows keep each row's own match in its top ten, the formal audit's convention.}\label{tab:replay}
\begin{tabular}{lrrrrr}
\toprule
query form & slots changed & queries changed & widest gap & certified, all & certified, changed \\
\midrule
float32 rows & 51 & 51 & $2.7\times 10^{-5}$ & 54.1\% & 0 \\
float16 roundings & 50 & 50 & $2.3\times 10^{-5}$ & 63.5\% & 0 \\
\midrule
float32 rows, self kept & 45 & & $2.5\times 10^{-5}$ & 57.9\% & 0 \\
float16 roundings, self kept & 38 & & $2.5\times 10^{-5}$ & 66.3\% & 0 \\
\bottomrule
\end{tabular}
\end{table}

\paragraph{The certificate on live stores, and its cost.} Read-only, 300 queries each: the
certificate certified the top ten of 59\% of stored vectors asked as queries on the main store
and 37\% on the curated store, and 50\% and 30\% with a text query's own rounding added
($\betaQ \approx 2.1\times 10^{-4}$). The fixture test certifies at least 20 of 60 reads on a
400-vector store at 64 dimensions and checks each certified set against the exact cosines; the
control constructs a query and two vectors whose exact and float16 rankings disagree, and asserts
that the read changes and is refused. On fixture stores of 2,000 and 17,256 rows at 768
dimensions, best of twenty passes of 60 queries interleaved with \texttt{main}, \texttt{search}
measured within 4\% of \texttt{main}'s either side, the noise between runs, and
\texttt{search\_ranked} 0.6 to 4\% over \texttt{search}.

\paragraph{The backward-error hypothesis.} On Laplacians of cycles, paths and complete graphs of
3 to 60 vertices, whose spectra are known in closed form, the computed eigenvalues' errors reach
$3.3\,nu\norm{A}$ (the SVD of a 19-cycle): an eigenvalue error, a lower bound on the backward
error, consistent with the hypothesis $8nu\norm{A}$ and not a proof of it. The Hodge rank cut is
decisive, and counts $\beta_1$ exactly, on a filled triangle, $K_4$, two squares, two holes at a
vertex and cycles of 4 to 39 vertices; the pseudo-inverse cut on 40 random weighted graphs of 10
to 120 vertices with weights in $[0.785, 1]$; the control, two triangles joined by a bridge of
weight $5\times 10^{-12}$, fails the gap condition and loses the bridge's mode. The rise
threshold's condition holds on at least 20 of 40 random graphs beside a second component and
fails on a path of forty notes with weight 0.9, where $g \approx 0.005$ and the bound passes
$10^{-9}$ while the measured movement stays below $10^{-3}$ of it.

\paragraph{The spine.} On synthetic spines of 2,801, 5,601, 11,201 and 22,401 events, the
real-store proper-time test body went from 1.97, 8.04, 36.8 and 166 s to 0.01, 0.01, 0.02 and
0.04 s; the first call per stratum still pays the closure.

\paragraph{The write-path races.} The compare-and-set test forces the interleaving with two handles
on one file, the second claiming between the first's select and update: the old code returns the
same job to both, the new claims it once, and with two jobs the loser takes the other row. The
pull request reports a further probe of six processes over 600 jobs, each claimed once; that probe
is not a committed test. The code lane's unlocked append was judged to race nothing under the
one-writer model, after four runs of three processes (two updating disjoint rows, one appending,
120 commits a run) with no error and no lost write.

\paragraph{Lineage and diamonds.} On the live ledger a path's chain members and its ledger
versions measured equal, 1,663 to 1,663, which prompted \cref{thm:lineage}. On the live versions
store, read-only, 29 identities hold one version each, so none forks or merges, and the
homotopy-coherent case stays deferred.

\subsection{The audits as evaluation}

Two audits ran over the branch before integration, one formal and one of its security. The formal
audit's findings are recorded in the commits that closed them; \cref{tab:audit} groups them by
class. The security audit's findings are not itemized on the branch; the commits that answer
write-path exposure (the refusals beside a live writer, the supervisor lock, the embed charge
bounded to a role's classes) are in the implementation section, each with a test. Every finding of
both audits was closed before integration.

\begin{table}[htbp]
\centering\footnotesize
\caption{The formal audit's findings, by class, and how each was closed.}\label{tab:audit}
\begin{tabular}{p{0.17\textwidth}p{0.38\textwidth}p{0.37\textwidth}}
\toprule
class & finding & closed by \\
\midrule
information flow & the oracle had a second caller, a script's embed charge, which told a role whether an atom exists in a class it cannot read & the charge reads only the role's classes (\texttt{known\_in}, \cref{thm:known}); a test lands an atom only one stratum holds and checks the charge \\
crash convergence & a delete that dies after its retire and before its tombstone, then the same bytes back, is skipped by UNCHANGED for ever & UNCHANGED requires current rows (\texttt{holds\_current}); the probe is the test, in the vault lane and the ingress writer \\
reproducibility & the design note's replay figures came from a NumPy reading, not the engine & \texttt{eval/float16\_replay.py} through the engine's kernels; the figures corrected \\
hypothesis scope & SPC-3's $\varepsilon$ covered the decomposition only, not forming the float64 Laplacian & $\varepsilon = 9nu\norm{L}$; the rise bound rederived with Wedin's $\mu$; a tautological test replaced by the hypothesis' consequence \\
test adequacy & no test held $\betaQ$: dropping the query term left every test green & a query one coordinate above a float16 tie, with the engine's distance as witness \\
statement precision & the status block claimed Result 4 holds, where TRN-15 checks the strict inclusion and one direction & the wording narrowed to what is checked \\
\bottomrule
\end{tabular}
\end{table}

\section{Limitations}

\paragraph{Synthetic fixtures, and figures reported but not committed.} The spectral-gap
fixtures (closed-form spectra, random weighted graphs with weights in $[0.785, 1]$, a 40-vertex
path), the certificate's fixture stores (2,000 and 17,256 rows; 400 vectors at 64 dimensions) and
the spine timings are synthetic. The replay is over an archived store of real rows, but one
corpus. The six-process claim probe and the three-process append probe are reported in the pull
request and commit messages; the branch holds no script for either. The live-store certification
rates and the 29-identity count were read once and are not re-measured by any test.

\paragraph{Hypotheses the code cannot check.} The float32 arithmetic of the search engine's
distance kernel is a stated bound $\gamma$, measured on fixtures and not modelled; the norm floor
$\nu = 1/2$ is a hypothesis on the corpus. NumPy's backward stability with constant $8n$ is the
LAPACK guide's statement, and the fixture measurements are consistent with it, not a proof.
Wedin's bound is cited and not formalized, so the rise threshold's condition is checked on
fixtures and fails on thin graphs, where the law is silent. The plane's noninterference still
rests on exact search with ties broken by atom id and on a content-addressed table, hypotheses
this branch does not discharge.

\paragraph{What the models leave out.} The plane model lands one fiber per transaction where the
code lands several; \lean{reset} and a relabel's place in a trace are outside it. The keyed door
model has no retire command, so the tombstone-last order is written to DOR-7's shape but not
cited from it. The lineage model has one path. The superconnection is on the nerve to its
2-skeleton, where every component of $\A^2$ lives; the homotopy-coherent case is not modelled,
whether a local isometry is necessary for Result 4 is not checked, and the push-forward reading
of its formulas is not separately checked.

\paragraph{The comparator locally.} \texttt{lake comparator} runs in a sandbox available on Linux,
where CI runs it. On the development machine \texttt{check.sh} runs an in-process comparison
instead, which compares statement types and definition values but not constructors, recursor
rules or axioms; a weakened statement and a changed one were each planted and caught, but the
local gate is weaker than CI's.

\section{Future work}

The keyed model should gain a retire command, so that the key-last law covers deletes as it covers
landings. A deterministic tie-break in the plane's search, one extra row per query and more only
at an exact tie, would remove the first of noninterference's hypotheses. The rise threshold could
scale with the conductance, which would make it provable on thin graphs. Whether sentence answers
should ask for the float16 certificate, and whether the existence oracle should be accepted,
bounded for its remaining caller or closed, are left open with the measurements they need. Result
3 should be amended with the no-restoration hypothesis and the correction's sign, and version
chains could be threaded over the first-parent line so that each edge is a commit's diff.

\section{Conclusion}

The three kinds of claim that remained at this stage of the ladder are now on it. The store's
write paths have models under which a migration, a landing, a relabel and an oracle are proved to
do what their prose said or shown to do otherwise; two did otherwise, and the code is written to
the model now. The float claims have a rounding model that meets the machine bit for bit and a
certificate the store computes, with a replay that locates every change float16 made in an
archived store inside the bound and outside the certificate. Three tolerances have the condition
under which they decide correctly, and one has moved to meet it. Two design results are
machine-checked with their errata.

\bibliographystyle{plainnat}
\bibliography{refs}
\end{document}
